% Section 6.3, from lectures/L06/appendix/c_c_abi.md.

\section{The C ABI}
\label{c6:sec:abi}

\subsection{The other party}
\label{c6:sec:party}
Since Chapter~1 you have been following a calling convention (\secref{c2:sec:contract}): arguments
in \code{r24} and \code{r22}, results in \code{r24}, \code{r18} to \code{r27} and \code{r30},
\code{r31} free to destroy, \code{r2} to \code{r17} and \code{r28}, \code{r29} to be given back.

Up to now that convention has had one party. Every caller of your subroutines was code you wrote,
so following the rules was a discipline rather than a requirement, and breaking them would only
ever have surprised you.

This section is the other party arriving. The C compiler follows the same convention, so a
function it generates can call a subroutine you wrote, and vice versa, with no glue and no
declaration beyond a prototype. \textbf{Everything you have been doing since Chapter~1 was for
this.}

\subsection{The convention, exactly}
\label{c6:sec:convention}
Figure~\ref{c2:fig:contract}, back in Chapter~2, has the registers coloured by role; this is the
rest of what it implies.

\textbf{Arguments go in \code{r25} downwards, in pairs, left to right.} The first argument takes
\code{r25:r24}, the second \code{r23:r22}, the third \code{r21:r20}, and so on down to \code{r8}.
An 8-bit argument takes only the lower register of its pair, so the first is \code{r24}, the second
\code{r22}, the third \code{r20}, and the odd-numbered ones are simply not used.

That is worth stating plainly because it is the rule people guess wrong: a byte argument does
\textbf{not} move up to fill the gap. Three byte arguments arrive in \code{r24}, \code{r22} and
\code{r20}.

\textbf{Results come back the same way.} A byte in \code{r24}, a 16-bit value in \code{r25:r24}, a
32-bit value in \code{r25:r22}.

\textbf{Anything larger than fits in the registers goes on the stack}, and so does any argument
past the last pair. Allocation runs from \code{r25} down to \code{r8}, which is eighteen registers
and therefore \textbf{nine pairs}: \code{r25:r24}, \code{r23:r22}, \code{r21:r20}, \code{r19:r18},
\code{r17:r16}, \code{r15:r14}, \code{r13:r12}, \code{r11:r10}, \code{r9:r8}. A byte argument takes
a pair and wastes the odd half rather than closing the gap, so six \code{uint8_t} arguments land in
\code{r24}, \code{r22}, \code{r20}, \code{r18}, \code{r16} and \code{r14}. This book never comes
close; a kernel's \code{task_init} does.

\textbf{\code{r1} is assumed to hold zero.} The compiler relies on it constantly and never checks.
A subroutine that uses \code{r1} as scratch and does not clear it before returning breaks C code
that had nothing to do with it, at a distance, in a way that looks like a compiler bug. If you use
\code{mul}, whose result lands in \code{r1:r0}, clear \code{r1} afterwards (\secref{c4:sec:stride}).

\textbf{\code{r0} is scratch for everybody} and nobody expects it preserved.

\subsection{Calling your assembly from C}
\label{c6:sec:asmfromc}
Two things are needed, and neither is a wrapper.

On the assembly side, the subroutine has to be visible to a linker. Everywhere else in this book
that is free, because \code{avra} makes every label global (\secref{c1:sec:avra}). \textbf{Here it
is not}, and that is the first of two reasons this chapter is the one place the book uses a second
assembler.

On the C side, a prototype whose types match what the subroutine actually expects:

\begin{ccode}
extern uint8_t shift_bits(uint8_t shifts);
\end{ccode}

That is all. The compiler puts the argument in \code{r24} because the prototype says it is a byte,
and reads the result from \code{r24} for the same reason, and your subroutine has been doing
exactly that since Chapter~1.

\textbf{The prototype is a promise nothing checks.} Declare the argument as \code{uint16_t} and the
compiler will pass it in \code{r25:r24} and your subroutine will read \code{r24} and ignore
\code{r25}. That gives the same answers as before, because C would have cut the argument to its low
byte anyway, and it is the dangerous kind of wrong: nothing changes, so nothing tells you. Declare
the result as \code{uint16_t} too and it stops being harmless. The caller reads \code{r25:r24},
\code{r25} still holds the argument's high byte because your subroutine never touched it, and
\code{shift_bits(300)} comes back as 256. There is no header, no name mangling and no link-time
type check; the linker matches a name to a name.

\subsubsection{The one place this book uses two assemblers}
\code{avr-gcc} drives GNU \code{as}, and GNU \code{as} does not read AVRASM2. It cannot assemble
\file{drivers/source/utils.asm}, and no flag makes it: the two are different languages describing
the same machine. \code{avra}, for its part, has no linker and cannot combine anything with a
compiled object file.

So for this section, and for nothing else in the book, you translate your own \code{shift_bits}
into GNU \code{as} and save it as \file{drivers/source/utils_gnu.S}. The instructions do not
change; only the directives around them do:

\begin{booktable}{@{}p{0.47\linewidth}L@{}}
  \thead{Your \file{utils.asm}} & \thead{\file{utils_gnu.S}} \\
  \midrule
  \code{.include "m328Pdef.inc"} & \texttt{\#include <avr/io.h>} \\
  \emph{(nothing; \code{avra} exports every label)} & \code{.global shift_bits} \\
  \code{.equ NAME = value} & \code{.equ NAME, value} \\
  \code{out PORTB, r16} & \code{out _SFR_IO_ADDR(PORTB), r16} \\
  \code{sts PORTB + 0x20, r16} & \code{sts PORTB, r16} \\
  \code{.org 0x0006} (words) & \code{.org 0x000C} (bytes) \\
\end{booktable}

\code{shift_bits} touches no port, defines no constant and sits in no vector slot, so in practice
the translation is: change the include, add one \code{.global} line, and copy the body across
unchanged. \textbf{The file must be \file{.S} and not \file{.s}}, because the capital is what runs
the preprocessor that makes \texttt{\#include} work.

\subsubsection{The same program, twice}
Here is a program that exercises every row of that table at once. It is not one of the book's
exercises; it exists to be read side by side. In AVRASM2, which is what you have been writing:

\begin{asm}
.include "m328Pdef.inc"

.equ LED_BIT = 5

.cseg
.org 0x0000
    rjmp reset
.org PCI0addr                   ; 0x0006, a WORD address
    rjmp on_pin_change

.org 0x0034                     ; past the table: 26 vectors, 2 words each
reset:                          ; global already; avra exports every label
    ldi r16, (1 << LED_BIT)
    out DDRB, r16               ; DDRB is 0x04, an I/O address
    sts TCCR1B, r16             ; 0x81, already a data space address
loop:
    rjmp loop

on_pin_change:
    in  r16, PINB               ; PINB is 0x03, an I/O address
    reti
\end{asm}

And the same program for \code{avr-gcc}, which is the dialect nearly all published AVR assembly is
written in:

\begin{gnuasm}
#include <avr/io.h>

.equ LED_BIT, 5                 /* comma, not = */

.section .text
.org 0x0000
    rjmp reset
.org 0x000C                     /* 0x0006 doubled: bytes, not words */
    rjmp on_pin_change

.org 0x0068                     /* 0x0034 doubled, for the same reason */
.global reset                   /* or nothing outside can call it */
reset:
    ldi r16, (1 << LED_BIT)
    out _SFR_IO_ADDR(DDRB), r16 /* 0x24 here; the macro subtracts 0x20 */
    sts TCCR1B, r16             /* unchanged: 0x81 in both dialects */
loop:
    rjmp loop

.global on_pin_change
on_pin_change:
    in  r16, _SFR_IO_ADDR(PINB)
    reti
\end{gnuasm}

\textbf{Every instruction is identical, and so is every byte.} Assemble both and disassemble the
results:

\begin{console}
   0:   33 c0           rjmp    .+102           ; to 0x68
   c:   32 c0           rjmp    .+100           ; to 0x72
        ...                                     ; the rest of the vector table
  68:   00 e2           ldi     r16, 0x20
  6a:   04 b9           out     0x04, r16
  6c:   00 93 81 00     sts     0x0081, r16
  70:   ff cf           rjmp    .-2
  72:   03 b1           in      r16, 0x03
  74:   18 95           reti
\end{console}

\textbf{The \code{.org 0x0034} is not padding.} Without it \code{reset:} starts at word \code{0x07},
which is the second word of PCINT0's slot, and the body runs straight through the slots for
PCINT1, PCINT2 and the watchdog. The program still works right up to the moment one of those
interrupts is enabled, and then the vector jumps into the middle of your setup code. Twenty-six
vectors of two words each end at word \code{0x33}, so \code{0x34} is the first word that is yours
(\secref{c3:sec:table}).

That listing came out of both files. Nothing in the two columns of the table is a difference of
\emph{machine}; every one of them is a difference of \emph{notation}, and the assembler resolves all
of it before a single byte is emitted. \code{out DDRB, r16} and \code{out _SFR_IO_ADDR(DDRB), r16}
are two ways of writing \code{b9 04}, and the chip has never heard of either.

Three of those rows are worth dwelling on:
\begin{itemize}
  \item \textbf{\code{.org PCI0addr} against \code{.org 0x000C}.} The same vector slot. AVRASM2
    counts words as the datasheet does, GNU \code{as} counts bytes, and the factor of two between
    them is the single most productive source of wrong-looking interrupt behaviour when code moves
    between the dialects (\secref{c3:sec:table}).
  \item \textbf{\code{DDRB} against \code{_SFR_IO_ADDR(DDRB)}.} Also the same register.
    \file{m328Pdef.inc} calls it \code{0x04} and \file{<avr/io.h>} calls it \code{0x24}, and the
    macro subtracts the \code{0x20} back off so that \code{out} can reach it
    (\secref{c1:sec:iowindow}).
  \item \textbf{\code{sts TCCR1B, r16}, which did not change at all.} \code{TCCR1B} lives above the
    I/O window, so both dialects give it the same data space address and neither can reach it with
    \code{out}. The registers that need translating are exactly the ones \code{out} can touch.
\end{itemize}

Doing this once by hand is the point. That is the whole of the difference between the dialect this
book writes and the dialect most published AVR assembly is written in, and you now have a reason to
read it properly (\secref{c1:sec:studio}).

\subsubsection{Building the two together}
Hand both files to \code{avr-gcc}, which recognises \file{.S} and \file{.c} and does the right thing
with each:

\begin{shell}
avr-gcc -mmcu=atmega328p -Os \
        drivers/app/main.c drivers/source/utils_gnu.S \
        -o drivers/build/mixed.elf
\end{shell}

Note that this is a \emph{different} build from the rest of the book, which assembled one unit
with \code{avra} and never involved a linker at all (\secref{c1:sec:avra}). A C program needs the
startup code: it is what sets up the stack pointer, clears \code{.bss}, sets \code{r1} to zero, and
calls \code{main}. All the things your \file{main.asm} had to do by hand.

\subsection{Calling C from your assembly}
\label{c6:sec:cfromasm}
The same convention read backwards. Put the arguments in \code{r25:r24} and \code{r23:r22},
\code{rcall} the function's name, and take the result from \code{r24}.

What changes is whose problem the registers are. \textbf{A C function will destroy \code{r18} to
\code{r27} and \code{r30}, \code{r31}}, because the convention says it may, and the compiler takes
full advantage. Assembly that calls C and expects \code{Z} to survive is relying on nothing.

And \textbf{the C function assumes \code{r1} is zero on entry}, so if your assembly has been using
it, put it back before the call and not merely before the return.

\subsection{Reading what the compiler produces}
\label{c6:sec:reading}
\code{avr-gcc -S} writes assembly instead of an object file, and it is worth doing at least once
with your own code beside it.

\begin{shell}
avr-gcc -mmcu=atmega328p -Os -S drivers/app/main.c -o -
\end{shell}

Two things are usually surprising the first time.

\textbf{The compiler is better at some things than you expect.} A function whose last act is to
call another becomes a \code{jmp} rather than a \code{call} and a \code{ret}. On this part the
compiler writes \code{call} and \code{jmp}, because 32\,KB of flash is further than \code{rcall}
and \code{rjmp} reach: the \code{call} and \code{ret} it removes cost eight cycles between them and
the \code{jmp} that replaces them costs three, so the saving is five cycles and two bytes of stack.
That is the tail call from Chapter~3's driver specification (\secref{c3:sec:irq}), applied
automatically and everywhere.

\textbf{And worse at others.} It does not know that a value fits in a byte when the C types say
\code{int}, it will not use a register you know is free, and it cannot make the arithmetic
decisions you made in Chapter~5. This chapter's cross-check (Exercise~\ref{c6:ex:crosscheck})
compares the compiler's \code{shift_bits} against yours, and the result is not the flattering one
you might expect in either direction.

\textbf{Read \code{-Os} output, not \code{-O0}.} Unoptimised AVR code spills everything to the
stack and is not representative of anything anybody ships.

\subsection{What the ABI does not give you}
\label{c6:sec:abinot}
\textbf{It does not check anything.} The prototype is a promise, the linker matches names, and a
mismatch is a program that runs and is wrong.

\textbf{It does not make your assembly portable.} The convention is \code{avr-gcc}'s. Another
compiler for the same chip may differ, and the same compiler for a different AVR family does differ
in the details of where large values go.

\textbf{And it does not make mixed code easy to reason about.} Every call is a boundary where one
side knows what it destroyed and the other knows what it needed, and neither can see the other.
That is the price of the arrangement, and the reason the convention is worth following from the
first subroutine you write rather than from the first one that C calls.

\subsection{Reading a C variable from assembly}
\label{c6:sec:variable}
Everything so far has been about \emph{calls}: C hands your subroutine registers, your subroutine
hands them back. There is a second half to the boundary and this book has not needed it once,
because every driver here is handed the address of everything it touches. \code{led_on} gets a
pointer in \code{r25:r24}; it never looks anything up by name.

Code that shares \emph{state} with C rather than being called by it does look things up by name,
and it is worth doing on paper before you need it.

\needspace{6\baselineskip}
Given, in some C file:

\begin{ccode}
volatile uint8_t  ticks;
volatile uint16_t deadline;
\end{ccode}

the assembly is:

\begin{gnuasm}
    lds  r24, ticks             ; one byte, one instruction
    inc  r24
    sts  ticks, r24

    lds  r24, deadline          ; two bytes, low first, two instructions
    lds  r25, deadline + 1
\end{gnuasm}

Three things in that are worth stating outright.

\textbf{A C global is a data space address, and the linker supplies it.} \code{ticks} is a name
for a byte in SRAM, somewhere above \code{0x0100}, chosen when the program was linked. That is why
it is \code{lds} and \code{sts} and never \code{in} and \code{out}: \code{in} and \code{out} reach
the first 64 I/O registers and nothing else, and no C variable ever lives there
(\secref{c1:sec:iowindow}). The mistake reads well and assembles cleanly, and the linker only
warns: it cuts the address down to the six bits \code{in} has room for and builds the program
anyway, so \code{in r24, ticks} reads whichever I/O register that lands on and never your variable.

\textbf{A 16-bit variable is two loads and the low byte is first.} Which means the same tearing
Chapter~3 warned about (\secref{c3:sec:shared}), on the same terms: two instructions with a gap,
and an interrupt that writes \code{deadline} can land in it.

\textbf{A pointer costs one more load than you think.} If C declares \code{struct thing* current},
then \code{current} is a \emph{pointer variable}: two bytes in SRAM holding the address of the
structure. Reading the structure's first field is two loads and then a third:

\begin{gnuasm}
    lds  r30, current           ; Z = the pointer's value, not its address
    lds  r31, current + 1
    ld   r24, Z                 ; and now the field
\end{gnuasm}

Loading \code{Z} with the address \emph{of} \code{current} and dereferencing that gives you the
pointer's own bytes, interpreted as data, which is a plausible-looking small number rather than an
error. It is the single most common way this boundary is got wrong, and the reason is that C hides
the difference: \code{current->field} is one arrow and two dereferences.

\textbf{There is no AVRASM2 half of this section}, and that is the point rather than an omission.
\code{avra} has no linker (\secref{c1:sec:avra}), so there is nothing to supply an address that
some other file chose; a name in a \file{.asm} file has to have been defined in the unit being
assembled. Sharing a variable with C is a thing you can only do in a build that links, which is
this chapter's build and no earlier one.

\textbf{Where this goes next.} A kernel is exactly the code that shares state with C rather than
being called by it. Its context switch is assembly, its scheduler is C, and the two meet across a
single \code{struct tcb*} that the C half assigns and the assembly half reads, dereferences and
writes back: the three instructions above, in that order, plus the stack pointer of
\secref{c4:sec:loadsp}.
